% formal/f64-orientation-sign-filters -- f64 error bounds for 4D orientation predicates
%
% Labels defined here:
%   def:f64-orient-gamma,
%   lem:f64-orient-product-error,
%   lem:f64-orient-det-error,
%   prop:f64-orient4-sign-filter,
%   alg:f64-origin-interior-triple-filter
%
% Status: agent-written proof note, unapproved by Jorn. The proof covers exact
% f64 inputs interpreted as real numbers. It does not justify using rounded f64
% values as proxies for separate exact rational inputs.

\providecommand{\fl}{\operatorname{fl}}

\subsection{f64 orientation sign filters}

The primitive covered here is
\[
  \texttt{orient4\_sign\_f64}(M)
  \in \{\textsc{positive},\textsc{negative},\textsc{indeterminate}\}
\]
for a finite f64 matrix $M\in\mathbb{F}^{4\times4}$, interpreted as an exact
real matrix.  The certified outputs mean
\[
  \textsc{positive}\Rightarrow \det(M)>0,\qquad
  \textsc{negative}\Rightarrow \det(M)<0.
\]
\textsc{Indeterminate} makes no claim.

\begin{definition}[Floating point error notation]\label{def:f64-orient-gamma}
Let $u=2^{-53}$ and $\gamma_k=ku/(1-ku)$ for $ku<1$.
We use the standard model
\[
  \fl(x\circ y)=(x\circ y)(1+\delta),\qquad |\delta|\le u,
\]
for $\circ\in\{+,-,\times\}$ only when the operation returns a finite normal f64
number or exact zero.  The Rust diagnostic returns \textsc{indeterminate} if an
input or intermediate value is non-finite or a nonzero subnormal.  If a
multiplication of two nonzero values rounds to zero, it also returns
\textsc{indeterminate}.  If an addition rounds to zero, the Rust code accepts
it only when the two f64 summands are exact negatives, so the real result of
that f64 addition is also zero.
\end{definition}

\begin{lemma}[Four-factor product]\label{lem:f64-orient-product-error}
Let $p=x_1x_2x_3x_4$ and compute
\[
  \hat p=\fl(\fl(\fl(x_1x_2)x_3)x_4).
\]
Under Definition~\ref{def:f64-orient-gamma},
\[
  \hat p=p(1+\theta_3),\qquad |\theta_3|\le\gamma_3.
\]
\end{lemma}

\begin{proof}
Compose the three multiplication roundoff factors and use the standard
$\gamma$-calculus bound.
\end{proof}

\begin{lemma}[Leibniz determinant error]\label{lem:f64-orient-det-error}
Let
\[
  \det(M)=\sum_{\pi\in S_4}\operatorname{sgn}(\pi)
    \prod_{i=0}^3 M_{i,\pi(i)}
\]
and let $\widehat{\det}$ be computed by evaluating the 24 four-factor products
as in Lemma~\ref{lem:f64-orient-product-error}, applying the signs exactly, and
summing the signed terms sequentially.  Let
\[
  P(M)=\sum_{\pi\in S_4}\prod_{i=0}^3 |M_{i,\pi(i)}|.
\]
Then
\[
  |\det(M)-\widehat{\det}|\le\gamma_{26}P(M).
\]
\end{lemma}

\begin{proof}
Each product contributes at most $\gamma_3$ relative error.  Sequentially
summing 24 signed terms contributes at most $\gamma_{23}$.  Composition gives
$\gamma_{3+23}=\gamma_{26}$.  Taking absolute values over the 24 exact product
magnitudes gives the stated bound; cancellation in the determinant does not
affect the bound because $P(M)$ is an absolute sum.
\end{proof}

\begin{proposition}[Computable sign filter]\label{prop:f64-orient4-sign-filter}
The Rust implementation also computes $\widehat P$, the sequential f64 sum of
the 24 rounded absolute products.  It then sets
\[
  \widehat E =
    \max\{\,\operatorname{nextUp}(\fl(\fl(16\gamma_{26})\widehat P)),
      \,\mathrm{minNormal}\,\}
\]
when $\widehat P>0$, and $\widehat E=0$ when $\widehat P=0$.  If the computation
of $\widehat E$ overflows or is NaN, including the final upward step, the result
is \textsc{indeterminate}.

For all non-indeterminate executions,
\[
  |\det(M)-\widehat{\det}| \le \widehat E.
\]
Therefore the decisions
\[
\begin{array}{rcl}
  \widehat{\det}>\widehat E &\Rightarrow& \textsc{positive},\\
  \widehat{\det}<-\widehat E &\Rightarrow& \textsc{negative},\\
  |\widehat{\det}|\le \widehat E &\Rightarrow& \textsc{indeterminate}
\end{array}
\]
are sound.
\end{proposition}

\begin{proof}
By Lemma~\ref{lem:f64-orient-product-error}, every exact absolute product is at
most its rounded product divided by $1-\gamma_3$.  For the absolute-product
sum, all summands are nonnegative.  In a non-indeterminate execution each
positive intermediate sum is normal; a zero sum can occur only when all rounded
absolute products seen so far are zero.  Therefore the standard summation bound
applies to the nonzero case, and the zero case has $P(M)=\widehat P=0$ because
nonzero products are not allowed to round to zero.  In the nonzero case, the
exact sum of the rounded absolute products is at most
$\widehat P/(1-\gamma_{23})$.  Hence
\[
  P(M)\le \frac{\widehat P}{(1-\gamma_3)(1-\gamma_{23})}.
\]
The computed value of $\gamma_{26}$ is formed from $26u/(1-26u)$ by one
multiplication, one subtraction, and one division.  In a non-indeterminate
execution it is therefore at least $(1-\gamma_3)\gamma_{26}$.  The two
multiplications forming $16\widehat\gamma_{26}\widehat P$ introduce at worst
another factor $(1-u)^2$ before the final rounded normal value is formed.
The implementation then applies one upward successor step, so the returned
normal bound is at least
\[
  16(1-\gamma_3)(1-u)^2\gamma_{26}\widehat P .
\]
It suffices to compare this with the previous bound for $P(M)$.  For
$u=2^{-53}$,
\[
  16(1-\gamma_3)^2(1-u)^2(1-\gamma_{23}) > 15,
\]
because $\gamma_3<4u$ and $\gamma_{23}<24u$.  Consequently the factor 16
dominates the use of $\widehat P$ in place of $P(M)$ and the final ordinary f64
roundings; the upward successor is extra padding and is not needed for the
displayed inequality.  If the rounded product for the bound is zero or subnormal, the
implementation returns $\mathrm{minNormal}$; in that branch the mathematical
product above is below the normal range, so the replacement is still
conservative.  Thus $\widehat E\ge\gamma_{26}P(M)$.  The sign implications
follow from
$\det(M)\in[\widehat{\det}-\widehat E,\widehat{\det}+\widehat E]$.
\end{proof}

\begin{algorithm}[Origin-interior triple-normal f64 filter]
\label{alg:f64-origin-interior-triple-filter}
For each triple $(i,j,k)$ and each non-triple point $\ell$, classify
\[
  \det(a_i,a_j,a_k,a_\ell),
\]
which is the sign of $a_\ell$ against the triple normal, up to the fixed
cross-product sign convention.

For a fixed triple:
\begin{itemize}
\item if some non-triple point is certified positive and some non-triple point
  is certified negative, the triple is certified nonseparating;
\item if every non-triple point is certified positive, the triple certifies
  $0\notin\operatorname{int}\operatorname{conv}\{a_1,\ldots,a_N\}$;
\item if every non-triple point is certified negative, the triple certifies
  $0\notin\operatorname{int}\operatorname{conv}\{a_1,\ldots,a_N\}$;
\item otherwise the triple is indeterminate.
\end{itemize}
The whole filter returns certified \textsc{false} as soon as a separating
triple is certified, certified \textsc{true} only when every triple is
certified nonseparating, and \textsc{indeterminate} otherwise.
\end{algorithm}

\begin{proof}[Soundness]
Every non-indeterminate sign is exact by
Proposition~\ref{prop:f64-orient4-sign-filter}.  A triple with certified signs
on both sides is not separating.  A triple with all non-triple points certified
strictly on the same side gives a nonzero normal whose hyperplane through the
origin supports all points on one side, so the origin is not interior.  All
other cases return \textsc{indeterminate}.

It remains to justify the certified \textsc{true} case.  Suppose, contrapositively,
that $0$ is not in the interior of $K=\operatorname{conv}\{a_1,\ldots,a_N\}$.
By the finite-dimensional separating hyperplane theorem, the cone
\[
  C=\{n\in\mathbb R^4:\langle n,a_i\rangle\ge 0\text{ for all }i\}
\]
contains a nonzero vector.  If the input has rank less than four, every
$4\times4$ determinant formed from input rows is zero, so every triple-normal
sign test is indeterminate and the algorithm cannot return \textsc{true}.

Assume therefore that the input points span $\mathbb R^4$.  Then
$C\cap(-C)=\{0\}$, so $C$ is a pointed polyhedral cone.  Choose an extreme ray
$\mathbb R_{\ge0}n$ of $C$.  For a pointed cone in $\mathbb R^4$, the active
constraints at an extreme ray span a codimension-one subspace: otherwise their
orthogonal complement inside $\mathbb R^4$ would have dimension at least two,
and the ray direction could be perturbed in two opposite directions while
remaining in all active constraints, contradicting extremality.  Hence there
are three linearly independent input points
$a_i,a_j,a_k$ with
\[
  \langle n,a_i\rangle=\langle n,a_j\rangle=\langle n,a_k\rangle=0.
\]
The corresponding triple normal is a nonzero scalar multiple of $n$.  Since
$n\in C$, all remaining input points lie in the same closed halfspace for that
triple normal.  Thus their exact determinant signs cannot include both a
strict positive and a strict negative sign.  The algorithm can certify a triple
as nonseparating only by certifying strict signs on both sides, so this triple
is not certified nonseparating.  Therefore a run that certifies every triple as
nonseparating has no separating hyperplane through the origin, and so
$0\in\operatorname{int}K$.
\end{proof}
